\chapter{Discussion}
\label{ch:discussion}

Chapter~\ref{ch:exp} reported what the hardware trials and the human-subject study measured; this chapter reads those results against the hypotheses of Section~\ref{sec:rqs}, weighs the alternative explanations the design does not rule out, and states the threats to validity that qualify every claim above.

\section{Findings against expectations}
\label{sec:discussion-expectations}

\textbf{H1} is rejected in the shape it predicted, although the interaction it predicted is real on one of its two named measures. The mode--assistance interaction on task time is significant (Table~\ref{tab:assumptions}), so the two mappings do respond differently to the assistance channel, but not as the reversal H1 stated: force guidance alone is the slowest setting under \emph{both} mappings, not only under joystick, so CF never leads the clutch mapping (Section~\ref{sec:resperf}); the corresponding interaction on path length, where the same clutch-side reasoning applies, is not significant ($p=0.090$). The clutch-side mechanism did not survive contact with the operators either: the pull toward the goal was reported as a force to resist, not a depth cue, and within the clutch mapping it is blending, not guidance, that lowers both the rendered force and the number of re-indexing operations (Section~\ref{sec:reseffort}). The joystick-side half fares better in ranking --- JF is the single slowest of the six conditions and JFB the single quickest --- and its channel-congruence mechanism is corroborated by a measure H1 did not name: operator--policy agreement $\rho$ shows an interaction of its own, the gain from adding guidance to blending several times larger under joystick than under clutch (Section~\ref{sec:resauthority}). That finding supports the reasoning behind H1's joystick half without confirming the time prediction itself, and the two channels together produce no detectable objective or subjective gain of FB over B in either mapping.

\textbf{H2} is not supported on either of its named measures. No mode effect on task time survives the practice check of Section~\ref{sec:reslearning}; where the objective record separates the mappings at all it leans the other way, the joystick mapping yielding the shorter hand path and the more confident intent estimate, while the clutch mapping's faster hand speed is spent on longer paths cut into about three times as many separate driving segments (Section~\ref{sec:resperf}); and the subjective composite index $\bar{S}$ shows no mode difference either, $p=0.498$, with the closing preference split at thirteen to eleven (Section~\ref{sec:ressubj}). The one-to-one correspondence between hand and gripper was as intuitive as expected once in motion, but the re-indexing procedure around it was not: several participants found the clutch sequence confusing and occasionally pressed the wrong button, an observation made during the sessions rather than a recorded measure, and the time and effort this procedure costs is what the objective measures register in place of the predicted advantage.

\textbf{H3} is rejected on both named measures, and in the opposite direction to the one predicted. Both workload subscales fall, rather than rise, under blending: mental demand from $4.35$ under force guidance alone to $3.21$--$3.23$ under the two blended settings, and physical demand from $3.62$ to $2.62$--$2.90$, every contrast Holm-corrected significant (Table~\ref{tab:subjassist}). The same two blended settings also raise the authority-cost items a felt intervention was expected to threaten: the sense of control climbs from $4.54$ to $5.77$ and trust from $4.42$ to $5.81$, and both are ranked ahead of guidance alone by nearly every participant (Figure~\ref{fig:subjranks}). The disturbance the hypothesis anticipated was felt, but on the other channel: operators described the guidance force, not the blended reference, as working against them and as overriding their intention, and some resisted it outright, while the blending correction --- although less counterable in principle, since the only way to oppose it is to move away from the goal --- drew few complaints. The scenario-specific half of H3 remains untested by design: the two scenarios are averaged into one value per condition before any statistic is computed (Section~\ref{sec:discussion-threats}), so no claim here is made about whether the effect is sharper in the cluttered relay scenario than in the independent one.

\textbf{H4} is supported on every named measure and exceeds its own stated scope. Blending lowers path length by about a third, tracking slack by about a quarter, and filter active time by about a third relative to guidance alone, while the attained clearance does not move (Sections~\ref{sec:resperf} and~\ref{sec:ressafety}). Task time, on which the hypothesis predicted no direction, falls with the others by about a quarter, so the correction of the executed trajectory paid in speed as well as in the quality H4 named.

Two findings fall outside all four hypotheses. The two transparency measures, mean and peak command alteration (Section~\ref{sec:resauthority}), disagree in direction between the mappings, and this is read here as a genuine finding about how arbitration behaves differently under a continuous-velocity mapping and a re-indexed position mapping, not as a contradiction to be explained away. The null result on minimum clearance (Section~\ref{sec:ressafety}) is read as the barrier itself setting an unmoving floor, consistent with its design, and not as evidence that the six conditions are equally safe in every other respect.

\section{Alternative explanations}
\label{sec:discussion-alternatives}

Several results admit an explanation other than the one Section~\ref{sec:conclfindings} states as the leading reading. The deficit of force guidance alone could reflect how its saturation was calibrated (Section~\ref{sec:study}) --- absolute in the clutch mapping, deadzone-exit-calibrated in the joystick mapping --- rather than a property of haptic guidance as a channel; a different calibration could narrow or widen the gap to blending. Blending's advantage on the subjective measures could be partly attributable to the visible reference marker, the colour cue that turns from orange to green as authority shifts (Section~\ref{sec:studyenv}), rather than to the arbitration law itself: the cue is present only in the blended conditions, so it is confounded with the channel under test. The joystick mapping's higher intent-estimator confidence (Section~\ref{sec:resauthority}) is a property of the estimator's own evidence stream --- a continuous velocity command supplies it steadily, a re-indexed position command interrupts it --- and need not reflect the operator's own intent being any clearer under one mapping than the other. The belief score's own competition-normalisation (Section~\ref{sec:belief}) can inflate apparent confidence once one goal candidate is even slightly ahead of its rivals, so the confidence measure may overstate how resolved the operator's intent actually was. The absence of a subjective mapping effect could be a ceiling or an order effect rather than genuine indifference, since ratings cluster toward the favourable end of the scale and mapping order was blocked rather than fully randomised (Section~\ref{sec:reslearning}). Finally, the monotonic worsening of operator effort across the session (Section~\ref{sec:reslearning}) is read here as fatigue, but a rising familiarity with how hard the handle can safely be pushed is not excluded by the data.

\section{Threats to validity}
\label{sec:discussion-threats}

\textbf{Internal.} Mapping order was blocked --- the three conditions of one mapping run before the three of the other, order balanced across participants --- rather than interleaved; the order check found the mode difference in composite score independent of which mapping came first ($p = 0.544$), but the same check failed for task time and for operator effort (Section~\ref{sec:reslearning}), so those two families carry a residual order risk the composite score does not. The two scenarios are averaged into one value per participant per condition, which cannot separate an effect specific to one scenario from a genuine main effect. No unassisted baseline was run, so no result here claims that assistance beats no assistance, only that the assistance channels differ from one another. The three-level assistance design (F, B, FB) does not realise the full $2\times2$ of the two channels, so channel main effects and their interaction are not separately estimable, only the F-versus-FB and B-versus-FB contrasts.

\textbf{Construct.} The rendered-force and re-indexing measures are not comparable across mappings, because the clutch tether and the joystick centring spring encode different physical quantities through the same actuator (Section~\ref{sec:reseffort}). The minimum clearance floor is a property of the barrier's own design margin, not of operator skill, which limits what that measure can be read to say about the assistance. Throughout, a non-significant test is treated as absence of evidence rather than evidence of equivalence, given $n=24$. The mode$\times$assistance interaction on peak command alteration flips significance between the corrected analysis of variance and the aligned rank transform (Section~\ref{sec:stats}); no claim in this thesis rests on that interaction.

\textbf{External.} The sample is twenty-four novices, predominantly male and reporting low prior familiarity with teleoperation devices (Section~\ref{sec:ressubj}), so the subjective preferences do not generalise to an experienced operator population without further evidence. The factorial comparison itself ran only in a matched simulation environment. The two hardware trials used a different, full-size haptic device and the robot's own live, camera-perceived world rather than the study's declared one (Section~\ref{sec:hwteleop}), so they corroborate the architecture's mechanisms without replicating the factorial result at hardware scale. The operator's picture throughout the study was a five-window display (Section~\ref{sec:studyapparatus}), a specific instrumentation choice that a different display configuration might not reproduce.

\textbf{Missing ablations.} Neither the alignment-based arbitration nor the adaptive dual-variable scheduling was isolated by ablation: no run compares the alignment law against a confidence- or distance-based arbitration rule, and no run compares the adaptive slack schedule against a fixed-weight baseline. Both mechanisms are implemented and exercised throughout Chapter~\ref{ch:exp}, but neither is shown here to outperform the alternative it was chosen over.

%───────────────────────────────────────────────────────────────────────────────
